\section{Compressione del modello per l'Edge}
\label{sec:compressione}

\subsection{Motivazione}

Un dispositivo Edge a risorse limitate come un Raspberry Pi dispone di
qualche gigabyte di RAM e di una CPU priva di acceleratore grafico dedicato.
Eseguire su un simile dispositivo un modello di Deep Learning nella sua
forma originale (33,44 MB, precisione a 32 bit per ogni parametro) è
possibile ma non ottimale in termini di occupazione di memoria e velocità di
risposta. La pipeline di compressione, implementata nello script standalone
\texttt{models/compress.py} (da eseguire manualmente e non richiamato
automaticamente dal ciclo principale), applica e confronta tre approcci
distinti.

\subsection{Pruning strutturato}

Il \emph{pruning} ("potatura") individua i filtri convoluzionali che
contribuiscono meno alla predizione finale (valutati tramite la norma L1 dei
rispettivi pesi) e li rimuove. Il progetto adotta lo \emph{structured
pruning}: a differenza del pruning non strutturato (che azzera singoli pesi
sparsi, senza ridurre effettivamente il numero di operazioni), il pruning
strutturato rimuove interi filtri, riducendo concretamente sia la dimensione
del modello sia il numero di operazioni da eseguire durante l'inferenza su
CPU. Nel progetto è stata applicata una sparsità del 30\% sui livelli
convoluzionali.

\subsection{Quantizzazione}

La \emph{quantizzazione} converte i pesi del modello dalla rappresentazione
a 32 bit in virgola mobile (float32) a una rappresentazione compatta a 8 bit
intera (INT8), riducendo la dimensione occupata di un fattore prossimo a 4
e velocizzando i calcoli su CPU prive di unità di calcolo dedicate al
floating point ad alta precisione. Il progetto ha sperimentato due percorsi
di quantizzazione:

\begin{itemize}
    \item \textbf{Dynamic quantization} in PyTorch: converte i pesi in INT8
    in anticipo, mentre le attivazioni (i dati che attraversano la rete
    durante l'inferenza) vengono quantizzate dinamicamente ad ogni
    esecuzione. È l'approccio più semplice e universalmente supportato, non
    richiede una fase di calibrazione.
    \item \textbf{Quantizzazione dinamica via ONNX Runtime}: il modello
    viene prima esportato in formato \textbf{ONNX} (\emph{Open Neural
    Network Exchange}), uno standard indipendente dal framework di
    provenienza che descrive sia il grafo computazionale sia i pesi in un
    unico file binario, eseguibile da un motore dedicato all'inferenza
    (\emph{ONNX Runtime}), tipicamente più efficiente di un framework di
    training generico come PyTorch quando utilizzato in sola inferenza. Il
    modello ONNX viene poi quantizzato con lo strumento
    \texttt{onnxruntime.quantization}.
\end{itemize}

La \emph{Post-Training Static Quantization} (PTQ statica), che quantizza
anche le attivazioni in anticipo tramite una fase di calibrazione su un
piccolo insieme di immagini di esempio, è stata valutata ma risultava
instabile sulla piattaforma di sviluppo (macOS con processore Apple
Silicon); si è pertanto optato per la quantizzazione dinamica, più robusta
e portabile.

\subsection{Le tre varianti compresse generate}

La pipeline produce tre varianti compresse distinte, ciascuna risultato di
una combinazione diversa delle tecniche precedenti, riassunte in
Tabella~\ref{tab:varianti-compresse}.

\begin{table}[h]
\centering
\small
\begin{tabular}{@{}p{2.6cm}p{4.4cm}p{5.6cm}@{}}
\toprule
\textbf{Variante} & \textbf{File generato} & \textbf{Tecnica applicata} \\
\midrule
\texttt{onnx} & \texttt{model\_quantized.onnx} & Sola quantizzazione dinamica INT8, tramite ONNX Runtime, sul modello originale non potato \\
\texttt{pruned\_quantized} & \texttt{model\_pruned\_quantized.pth} & Pruning strutturato (30\%) + quantizzazione dinamica INT8, in PyTorch \\
\texttt{pruned} & \texttt{model\_pruned.pth} & Solo pruning strutturato (30\%), pesi ancora in float32 \\
\bottomrule
\end{tabular}
\caption{Le tre varianti compresse prodotte dalla pipeline di \texttt{compress.py}.}
\label{tab:varianti-compresse}
\end{table}

Il motore di inferenza (\texttt{InferenceEngine}) può caricare una qualsiasi
di queste varianti in base alla configurazione
(\texttt{model.optimized\_variant} in \texttt{config.yaml}, oppure
selezionabile direttamente dalla dashboard), rendendo possibile un confronto
sperimentale diretto, i cui risultati sono riportati nella
Sezione~\ref{sec:risultati}.

\subsection{Problematiche tecniche incontrate e risolte}

Nel corso dello sviluppo sono emerse, e sono state diagnosticate e risolte,
alcune problematiche tecniche non banali, che vale la pena documentare
poiché rappresentano un aspetto significativo del lavoro di ingegnerizzazione
del progetto, oltre alla sola implementazione della pipeline "sul caso
ideale".

\subsubsection{Data leakage nella valutazione dell'accuratezza}

Nella prima versione del sistema, la fotocamera virtuale campionava
immagini dall'intero dataset PlantVillage-Tomato, incluso il sottoinsieme
utilizzato per l'addestramento della rete nel progetto originario. Questo
comportamento, noto in letteratura come \emph{data leakage}, produceva
un'accuratezza osservata prossima al 100\%, non rappresentativa della reale
capacità di generalizzazione del modello, poiché la CNN si trovava spesso a
classificare immagini già osservate durante il proprio addestramento. Il
problema è stato risolto scrivendo uno script dedicato
(\texttt{scripts/extract\_test\_set.py}) che ricostruisce esattamente lo
split casuale (\texttt{seed=42}, proporzione 80/20) utilizzato nel progetto
di training originale, isolando così il 20\% di immagini realmente escluse
dall'addestramento (2.906 immagini su un totale di 14.529), e configurando
la fotocamera virtuale affinché attinga esclusivamente da questo
sottoinsieme. Dopo la correzione, l'accuratezza osservata sul modello
originale si è attestata al 97,67\% su un campione di verifica di 300
immagini, valore coerente con quanto misurato nel progetto di training
originale.

\subsubsection{Fallimento silenzioso della quantizzazione ONNX}

Durante la verifica sperimentale delle diverse varianti compresse, si è
osservato che la variante \texttt{onnx} restituiva un'accuratezza
\emph{identica, bit per bit}, a quella del modello originale non
compresso -- un risultato sospetto, poiché una quantizzazione reale, per
quanto conservativa, introduce sempre una minima perdita di precisione
numerica. L'ispezione diretta del file \texttt{model\_quantized.onnx} ha
rivelato che esso era, di fatto, una copia identica del modello non
quantizzato. La causa risiedeva in un'eccezione sollevata internamente da
\texttt{onnxruntime.quantization.quantize\_dynamic()} durante l'esportazione
del grafo ONNX, generato con l'esportatore "dynamo" di PyTorch (diventato
predefinito nelle versioni recenti del framework): tale esportatore produce,
per l'operazione di \emph{flatten} presente nell'architettura
(\texttt{x.view(x.size(0), -1)}), un grafo con un'incoerenza nell'inferenza
automatica delle dimensioni dei tensori, che fa fallire il passo di
quantizzazione. Il codice della pipeline, intercettando genericamente
l'eccezione con un blocco \texttt{try/except}, ripiegava silenziosamente
sulla copia del modello non quantizzato, mascherando così il problema. La
correzione è consistita nel forzare l'utilizzo dell'esportatore ONNX
"legacy" (basato su TorchScript, tramite il parametro
\texttt{dynamo=False} di \texttt{torch.onnx.export}), risolvendo
l'incompatibilità e ottenendo una quantizzazione realmente funzionante
(dimensione del file ridotta da 33,44 MB a 8,37 MB, con accuratezza
sostanzialmente invariata).

\subsubsection{Backend di quantizzazione non inizializzato}

Un secondo problema, di natura simile, si presentava nel caricamento della
variante \texttt{pruned\_quantized}: il caricamento del checkpoint
PyTorch quantizzato falliva con l'errore \texttt{RuntimeError: Unknown
qengine}. La causa era l'assenza, nel motore di inferenza,
dell'impostazione esplicita del \emph{backend di quantizzazione}
(\texttt{torch.backends.quantized.engine}), un parametro globale di PyTorch
che deve essere impostato coerentemente sia in fase di creazione sia in
fase di caricamento di un modello quantizzato. Il problema è stato risolto
impostando esplicitamente il backend appropriato (\texttt{qnnpack} su
architetture ARM/Apple Silicon, \texttt{fbgemm} su x86) all'importazione del
modulo di inferenza.

\begin{notebox}
La documentazione di entrambe le correzioni non ha valore puramente
tecnico: dimostra un processo di verifica sperimentale rigoroso, in cui i
risultati numerici non sono stati accettati acriticamente ma confrontati e
messi in discussione quando apparivano statisticamente sospetti (accuratezza
identica fra baseline e modello quantizzato), portando all'identificazione
di problemi altrimenti silenti nella pipeline.
\end{notebox}
